% Part: first-order-logic
% Chapter: natural-deduction
% Section: provability-consistency

\documentclass[../../../include/open-logic-section]{subfiles}

\begin{document}

\iftag{FOL}
      {\olfileid{fol}{ntd}{prv}}
      {\olfileid{pl}{ntd}{prv}}

\olsection{\usetoken{S}{derivability} మరియు అవైరుధ్యం}

ఇప్పుడు \tetoken{వ్యుత్పాద్యత}{!!{derivability}} సంబంధపు అనేక ధర్మాలను స్థాపిస్తాం. అవి
స్వతంత్రంగా ఆసక్తికరమైనవి; అంతేకాక ప్రతి ఒక్కటీ సంపూర్ణతా సిద్ధాంతపు
నిరూపణలో ఒక పాత్ర పోషిస్తుంది.

\begin{prop}\ollabel{prop:provability-contr}
  $\Gamma \Proves !A$ అయి, $\Gamma \cup \{!A\}$ వైరుధ్యభరితమైతే,
  $\Gamma$ వైరుధ్యభరితమైనది.
\end{prop}

\begin{proof}
$\Gamma$ నుంచి $!A$ యొక్క \tetoken{వ్యుత్పత్తిని}{!!{derivation}} $\delta_1$ అనీ,
$\Gamma \cup \{!A\}$ నుంచి $\lfalse$ యొక్క \tetoken{వ్యుత్పత్తిని}{!!{derivation}}
$\delta_2$ అనీ అందాం. అప్పుడు కిందివిధంగా \tetoken{నిగమనం}{!!{derive}} చేయవచ్చు:
\begin{prooftree}
\AxiomC{$\Gamma, \Discharge{!A}{1}$}
\RightLabel{$\delta_2$}
\DeduceC{$\lfalse$}
\DischargeRule{\Intro{\lnot}}{1}
\UnaryInfC{$\lnot !A$}
\AxiomC{$\Gamma$}
\RightLabel{$\delta_1$}
\DeduceC{$!A$}
\RightLabel{\Elim{\lnot}}
\BinaryInfC{$\lfalse$}
\end{prooftree}
కొత్త \tetoken{వ్యుత్పత్తిలో}{!!{derivation}} పరికల్పన~$!A$ \tetoken{ఉపసంహరించిన}{!!{discharged}} అయింది; కాబట్టి
అది $\Gamma$ నుంచి ఒక \tetoken{వ్యుత్పత్తి}{!!a{derivation}}.
\end{proof}

\begin{prop}
\ollabel{prop:prov-incons}
$\Gamma \Proves !A$ అయితే, అప్పుడు మాత్రమే $\Gamma \cup \{\lnot !A\}$
వైరుధ్యభరితమైనది.
\end{prop}

\begin{proof}
మొదట $\Gamma \Proves !A$ అనుకుందాం; అంటే \tetoken{ఉపసంహరించని}{!!{undischarged}}
పరికల్పనలు~$\Gamma$ నుంచి $!A$ యొక్క ఒక \tetoken{వ్యుత్పత్తి}{!!a{derivation}}~$\delta_0$
ఉంది. కిందివిధంగా $\Gamma \cup \{\lnot !A\}$ నుంచి $\lfalse$ యొక్క
ఒక \tetoken{వ్యుత్పత్తిని}{!!a{derivation}} పొందుతాం:
\begin{prooftree}
  \AxiomC{$\lnot !A$}
  \AxiomC{$\Gamma$}
  \RightLabel{$\delta_0$}
  \DeduceC{$!A$}
  \RightLabel{\Elim{\lnot}}
  \BinaryInfC{$\lfalse$}
\end{prooftree}

ఇప్పుడు $\Gamma \cup \{\lnot !A\}$ వైరుధ్యభరితమని అనుకుందాం;
దానికి అనుగుణంగా $\Gamma \cup \{\lnot !A\}$లోని \tetoken{ఉపసంహరించని}{!!{undischarged}}
పరికల్పనల నుంచి $\lfalse$ యొక్క \tetoken{వ్యుత్పత్తిని}{!!{derivation}} $\delta_1$ అందాం.
$\FalseCl$ను ఉపయోగించి $\Gamma$ ఒక్కదాని నుంచి $!A$ యొక్క ఒక
\tetoken{వ్యుత్పత్తిని}{!!a{derivation}} ఇలా పొందుతాం:
\begin{prooftree}
  \AxiomC{$\Gamma, \Discharge{\lnot !A}{1}$}
  \RightLabel{$\delta_1$}
  \DeduceC{$\lfalse$}
  \RightLabel{\FalseCl}
  \DischargeRule{\FalseCl}{1}
  \UnaryInfC{$!A$}
\end{prooftree}
\end{proof}

\begin{prob}
$\Gamma \Proves \lnot !A$ అయితే, అప్పుడు మాత్రమే
$\Gamma \cup \{!A\}$ వైరుధ్యభరితమని నిరూపించండి.
\end{prob}

\begin{prop}\ollabel{prop:explicit-inc}
  $\Gamma \Proves !A$ అయి, $\lnot !A \in \Gamma$ అయితే $\Gamma$
  వైరుధ్యభరితమైనది.
\end{prop}

\begin{proof}
  $\Gamma \Proves !A$, $\lnot !A \in \Gamma$ అనుకుందాం. అప్పుడు
  $\Gamma$ నుంచి $!A$ యొక్క ఒక \tetoken{వ్యుత్పత్తి}{!!a{derivation}}~$\delta$ ఉంటుంది.
  $\Elim{\lnot}$ నియమాన్ని ఇలా సరళంగా వర్తింపజేయండి:
  \begin{prooftree}
    \AxiomC{$\lnot !A$}
    \AxiomC{$\Gamma$}
    \RightLabel{$\delta$}
    \DeduceC{$!A$}
    \RightLabel{\Elim{\lnot}}
    \BinaryInfC{$\lfalse$}
  \end{prooftree}
  $\lnot !A \in \Gamma$ కనుక అన్ని \tetoken{ఉపసంహరించని}{!!{undischarged}} పరికల్పనలూ
  $\Gamma$లోనే ఉన్నాయి; అందువల్ల $\Gamma \Proves \lfalse$.
\end{proof}

\begin{prop}\ollabel{prop:provability-exhaustive}
  $\Gamma \cup \{!A\}$, $\Gamma \cup \{\lnot !A\}$ రెండూ
  వైరుధ్యభరితమైతే, $\Gamma$ వైరుధ్యభరితమైనది.
\end{prop}

\begin{proof}
$\Gamma \cup \{ !A \}$ నుంచి $\lfalse$ను, $\Gamma \cup \{ \lnot !A \}$
నుంచి $\lfalse$ను ఇచ్చే \tetoken{వ్యుత్పత్తులు}{!!{derivation}s}~$\delta_1$, $\delta_2$ వరుసగా ఉన్నాయి.
అప్పుడు కిందివిధంగా \tetoken{నిగమనం}{!!{derive}} చేయవచ్చు:
\begin{prooftree}
\AxiomC{$\Gamma, \Discharge{\lnot !A}{2}$}
\RightLabel{$\delta_2$}
\DeduceC{$\lfalse$}
\DischargeRule{\Intro{\lnot}}{2}
\UnaryInfC{$\lnot \lnot !A$}
\AxiomC{$\Gamma, \Discharge{!A}{1}$}
\RightLabel{$\delta_1$}
\DeduceC{$\lfalse$}
\DischargeRule{\Intro{\lnot}}{1}
\UnaryInfC{$\lnot !A$}
\RightLabel{\Elim{\lnot}}
\BinaryInfC{$\lfalse$}
\end{prooftree}
$!A$, $\lnot !A$ పరికల్పనలు \tetoken{ఉపసంహరించిన}{!!{discharged}} అయ్యాయి కనుక ఇది
$\Gamma$ ఒక్కదాని నుంచి $\lfalse$ యొక్క ఒక \tetoken{వ్యుత్పత్తి}{!!a{derivation}}. అందువల్ల
$\Gamma$ వైరుధ్యభరితమైనది.
\end{proof}

\end{document}
